\documentclass[12pt]{article}

\usepackage{amsmath}
\usepackage{bm}

\begin{document}

pag55


\hfill

\textbf{3 Mathematical Introduction to Dynamical
Systems}

\hfill

No roote, no fruite

\hfill

Jeremiah Dyke, A worthy communicant,

\hfill

(1640)

\hfill

In this chapter, we present in the simplest possible manner a survey of some of the
fundamental mathematical concepts and tools which are required for the qualitative
analysis of the long-term behaviour of dynamical systems. A knowledge of the
theory of linear differential equations is a pre-requisite for the comprehension of
non-linear dynamics. The reader can find more detailed discussions in Chapters 5
and 6 of this book.

\hfill

\textbf{3.1 Linear Autonomous Systems}

\hfill

As we already mentioned in Chapter 2, we consider dynamical processes of a finite
dimension n which are differentiable and deterministic.

\hfill

In the simplest case, the system of differential equations (2.3.3) is linear with constant
coefficients

\[
\mathbf{\dot{x} = Lx }
\]

where \textbf{L} denotes the constant, non-singular ($n \times n$)-matrix of the coefficients. Using
this example, we now demonstrate which solutions can occur and which patterns
the relevant trajectories form in the corresponding phase space.

\hfill

We introduce the ansatz

\[
\mathbf{x} = e^{\lambda t} \mathbf{y}, \hspace{3mm} \mathbf{y} = \{y_1 y_2, ... y_n \}
\]

into eq. (3.1.1), thus obtaining a homogeneous linear system of equations

\[
[ \mathbf{L} - \lambda \mathbf{I}] \mathbf{y} = 0
\]

which only possesses non-trivial solutions if the system determinant vanishes

\[
P(\lambda) = \det[\mathbf{L} - \lambda \mathbf{I}] = 0
\]

or

\[
P(\lambda) = A_n \lambda^n + A_{n-1} \lambda^{n-1} + ... + a_0 = 0
\]

$P(\lambda)$ is a polynomial of the nth degree in $\lambda$ and is called characteristic or secular
equation. The zeroes of $P(\lambda)$ are the eigenvalues of \textbf{L}. A non-vanishing vector \textbf{y} 


which satisfies eq. (3.1.3) is called the eigenvector of L appertaining to the eigenvalue
$\lambda$. If $\lambda$ and y satisfy eq. (3.1.3), eq. (3.1.2) is a solution of eq. (3.1.1). For each
pair $\lambda_i,y_i$, we obtain in accordance with eq. (3.1.2) a solution of the form

\[
x_i = e^{\lambda_i t}y_i 
\]

If all n eigenvalues $\lambda_i$ are different, n linearly independent eigenvectors $y_i$ exist and
the general solution of eq. (3.1.1) can be expressed by the linear combination

\[
x(t) = \sum_{i=1}^n C_i e^{\lambda_i t} y_i
\]

with n integration constants Ci which are determined by the initial conditions $x(t_0)$
(Arnold, 1980; Braun, 1979).

\hfill

From a physical point of view, finding all the states of equilibrium $x_s$ is of interest.
These are stationary states of the system in which the dynamical process does
not undergo any change; they are thus characterised by $\dot{x}_s = 0$. If we observe the
corresponding point in the phase space, the vector $F(x_s)$ which defines the change of 


the trajectory vanishes here. For this reason, xs is thus also denoted singular point.
The decisive factor in characterising the stationary state is the behaviour of the
trajectories in the neighbourhood of the singularity. We speak of stable, unstable
or neutral equilibrium. If all the trajectories within certain neighbourhoods of xs are
captured, the singular point is asymptotically stable (sink); if, one the other hand,
all the trajectories that come close enough to xs are repelled, xs is asymptotically
unstable (source). On the basis of a system of two differential equations, we now
demonstrate the classification of singularities

\[
\mathbf{\dot{x} = Lx}, \hspace{5mm} \mathbf{L} = \begin{bmatrix}
L_{11} &  L_{12} \\
L_{21} & L_{22} \\
\end{bmatrix}, \hspace{5mm} \mathbf{x} \{x_1, x_2 \}
\]

with the origin $\mathbf{x}_s = 0$ as the singular point. The eigenvalues of the corresponding
eq. (3.1.3) serve as a basis to make the following six distinctions, illustrated in the
diagram fig. 3.1.1; here, we ignore the degenerate case of a zero eigenvalue.

\[
(A) \hspace{10mm}  \lambda_1 \neq \lambda_2
\]

Initially, we assume that \textbf{L} possesses different eigenvalues $\lambda_1 \neq \lambda_2$ and denote the
corresponding eigenvectors as $\mathbf{y}_1, \mathbf{y}_2$ which we can presume to be normalised without
loss of generality.

\hfill

We now introduce a new system of coordinates

\[
\mathbf{\bar{x}} = \{ \bar{x}_1 \bar{x}_2 \} 
\]

which is defined by the transformation

\[
\mathbf{x} = \mathbf{T \bar{x}} \text{ with } \mathbf{T} = [ \mathbf{y}_1 \mathbf{y}_2 ] 
\]

Thus, the eigenvectors $\mathbf{y}_1$, $\mathbf{y}_2$ are the basis vectors of the new coordinate system.
Substitution of the transformation eq.(3.1.10) into the given system eq.(3.1.8)
yields

\[
\mathbf{\dot{\bar{x}}} = \mathbf{D \bar{x}} \hspace{3mm} with \hspace{3mm} \mathbf{D = T^{-1}LT} 
\]

Due to

\[
\lambda_i y_i = \mathbf{L y_i}
\]a

\textbf{D} can be simplified as follows

\[
D = [ y_1 \ \ y_2 ]^{-1} L [y_1 \ \ y_2] = [y_1 \ \ y_2]^{-1} [y_1 \ \ y_2 ] [\lambda_1 \ \ \lambda_2]
\]
\[
= [\lambda_1 \lambda_2] 
\]

and is thus reduced to diagonal form. The transformed system equivalent to
eq. (3.1.8) now becomes

\[
\begin{matrix}
\dot{\bar{x}}_1 = \lambda_1 \bar{x}_1 \\
\dot{\bar{x}}_2 = \lambda_1 \bar{x}_2
\end{matrix}
\]

Here, we can distinguish for $\lambda_1 \neq \lambda_2$ a total of four possibilities (see fig. 3.1.1). If
the eigenvalues are real, the solution of eq. (3.1.13) according to eq. (3.1.7) is

\[
\bar{x}_1 = C_1 e^{\lambda_1 t}, \hspace{3mm} \bar{x}_2 = C_2 e^{\lambda_2 t}
\]

or, after eliminating the time parameter t,

\[
\bar{x}_2 = C x_1^{\lambda_2 / \lambda_1}
\]

We next differentiate between two cases having different signs (see fig. 3.1.1);

\[
(\textbf{A1}) \hspace{10mm} sign \lambda_1 = sign \lambda_2
\]
 
If the eigenvalues have the same sign, eq. (3.1.15) describes a family of parabolae
of the order $\lambda_2/ \lambda_1$ which have a common tangent at their origin (see fig. 3.1.2).
We then speak of a node. If the eigenvalues are negative, the node is stable, as can
be seen in the parameter representation of eq. (3.1.14) since all trajectories tend
towards the singular point for $t \to \infty$.

\[
(\textbf{A1}) \hspace{10mm} sign \lambda_1 = sign \lambda_2
\]

If $\lambda_1$ and $\lambda_2$ possess different signs, the solution is

\[
\bar{x}_2 \bar{x}_1^{|-\lambda_2 / \lambda_1|} = C
\]

i.e. the trajectories are hyperbolae. The corresponding singular point is then called
a saddle point or hyperbolic point (see fig. 3.1.3).

\hfill

If $\lambda_1, \lambda_2$ are complex conjugate, \textbf{L} can always be reduced by applying a (real)
linear coordinate transformation $\hat{T}$ to the simple form

\[
\dot{\hat{\mathbf{x}}} = \hat{L}x \text{ with } \hat{L} = \begin{bmatrix}
\alpha & -w \\
w & \alpha \\
\end{bmatrix}
\]

The eigenvectors $\mathbf{y}_1$ und $\mathbf{y}_2$ appertaining to the eigenvalues $\lambda_{1,2} = \alpha \pm  i w$ are also
complex conjugate

\[
\mathbf{y}_1 = \{ 1 - i \} \text{ and } \mathbf{y}_2 = \{ 1 i \} = \mathbf{y}_1^{*}
\]

and reduce in accordance with eq. (3.1.10) the transformation of $\hat{L}$ to a diagonal form; in accordance with standard usage, $\mathbf{y}_1^*$ denotes the complex conjugate vector
of $\mathbf{y}_1$. Thus, points $\hat{x}$ of the real phase space are transformed into points $\bar{x}$ of a two-dimensional
complex phase space $ \mathbf{C}^2$. The solution of eq. (3.1.13) can be expressed
immediately as

\[
\bar{x}_{1,2}= C_{1,2} e^{(\alpha \pm i w)t)} = C_{1,2}e^{\alpha t} (\cos wt \pm i \sin wt)
\]

where the complex constants $C_1$ and $C_2$ coincide with the initial values $\bar{x}_1(0)$ and
$\bar{x}_2$. It is now easy to see that $C_1$ and $C_2$ are complex conjugate. From the inverse
transformation of eq. (3.1.10) for an initial vector $\hat{\mathbf{s}}(0)$

\[
\bar{\mathbf{x}}(0) = \mathbf{T}^{-1} \hat{x}(0) = [\mathbf{y}_1 \ \ \mathbf{y}_1^*]^{-1} \hat{\mathbf{x}}(0) = \frac{1}{2} \begin{bmatrix}
\mathbf{y_1^{*t}} \\
\mathbf{y}_1^t
\end{bmatrix} \hat{x}(0)
\]

we obtain the relation $\bar{x}_2(0) = \bar{x}_1^*(0)$, i.e. $C_2 = C_1^*$. Thus, the solutions according
to eq. (3.1.17) are also complex conjugate, i.e. $\bar{x}_2(t) = \bar{x}^*_1(t)$ applies. The inverse
transformation of the initial system hence yields a real solution

\[
\hat{x}(t) = e^{\alpha t}\begin{bmatrix}
\cos wt & -\sin wt \\
\sin wt & \cos wt \\
\end{bmatrix} \hat{x}(0)
\]

We observe that the initial vector $\hat{x}(0)$ is rotated about an angle wt and stretched with the factor $e^{\alpha t}$.

\hfill

We once again distinguish between two cases (see fig. 3.1.1):

\[
(\textbf{A3}) \alpha = Re(\alpha_i) \neq 0
\]

Equation (3.1.18) suggests a representation in polar coordinates for which we set

\[
r = r_0 e^{\alpha t}, \varphi= wt + \varphi_0
\]

with $r_0 = | \hat{x}(0)|$ and $\varphi_0 = \tan^{-1}( \hat{x}_2(0)/ \hat{x}_1(0))$. This formulation describes a family
of logarithmic spirals. In accordance with an Anglo-Saxon convention, we use the
inverse function sign for arcus. If the real part of $\lambda_i$ is negative, i.e. $Re(\lambda_i) = \alpha < 0$,
we have a stable focus or spiral point; on the other hand, for $\alpha > 0$, the singular
point is unstable. Figure 3.1.4 shows the spirals in the original $x_1, x_2$ coordinate
system. They are affinely distorted by the inverse transformation.

\[
(\textbf{A4}) \alpha = Re(\lambda_i) = 0
\]

In the special case of $\alpha = 0$, the phase curves in the $\hat{x}_1, \hat{x}_2$ coordinate system
become concentric circles in accordance with eq. (3.1.18), the centre of which is the
singular point. We then speak of a centre. Figure 3.1.5 shows the trajectories in
the original $x_1, x_2$ coordinate system in which \textbf{L} possesses a general form. Here, we
observe concentric ellipses.

\[
(\textbf{B}) \lambda_1 = \lambda_2 = \lambda
\]

We now move on to a discussion of the case of identical eigenvalues, $\lambda_1 = \lambda_2 = \lambda$.
Here, we distinguish between two cases (fig. 3.1.1):

\hfill

(\textbf{B1}) Therearetwo linearly independent eigenvectors $matbf{y}_1, \mathbf{y}_2$ for $\lambda$.

In analogy to eqs. (3.1.11), (3.1.12), the system x˙ = Lx can be transformed into a
diagonal form, $\dot{\bar{x}} = \lambda \bar{x}_i$ (i = 1, 2), with the solutions

\[
\bar{x}_i = C_i e^{\lambda t} (i = 1, 2) 
\]

or

\[
\bar{x}_2 = C \bar{x}_1 
\]


We thus obtain a bundle of straight lines and speak of a dicritical node (see
fig. 3.1.6).

\hfill

In this case, a sufficient number of eigenvectors does not exist to span the whole
phase space; thus, \textbf{L} can no longer be transformed to diagonal form. However, in
matrix theory, it is shown that any non-diagonalisable (n × n) matrix can always
be transformed uniquely to the so-called Jordan normal form where each Jordan
block is associated with an eigenvector (Zurm¨uhl, 1964).

\hfill

For n = 2, there is thus a specific $\bar{x}_1, \bar{x}_2$ coordinate system in which the system of
differential equations (3.1.8) can be transformed to the Jordan normal form


\[
\dot{\bar{\mathbf{x}}} = \mathbf{\bar{L} \bar{x}} \text{ with } \bar{\mathbf{L}} = \begin{bmatrix}
\lambda & 1 \\
0 & \lambda  \\
\end{bmatrix}
\]

The corresponding coordinate transformation is

\hfill

(B2) There is only one linearly independent eigenvector \textbf{y} for $\lambda$.

\hfill

x = T ¯x with T = [y z] (3.1.23)
where a second vector z linearly independent of y must be established. This vector
z is denoted generalised eigenvector. Pre-multiplying
¯L
L = T
−1LT LT
with T yields
T¯LL = LT
Evaluating the matrix products, we find
T¯LL = [y z]
 
λ 1
0 λ
 
= [λy (y + λz)]
and
LT = L[y z] = [λy Lz]
Comparing the second columns and resolving for y yields
[L − λI]z = y




\end{document}